% Nouveau EGA V (cote 25): archive cote 25, site page 107 (PDF page 108).
% typescript. Draft compared: none.
% AI-transcribed by Claude Opus 5.5 (claude-opus-5-5, adaptive thinking, effort high), Message Batches API, batch ega25b (msgbatch_01A7n5YuXubUPNK1zsepi8RG), 2026-11-13; not yet checked by eye.
\scanpage{25/107}
\mspage{10}

Lorsque D $\geq 0$, cela implique $n \geq 0$, car si \inserted{\unsure{on avait|n'avait}} $n \unsure{<} 0$ toute section de
$\underline{L}$ sur X serait nulle. Grâce à 1.1. , si D est un diviseur relatif
$\geq 0$ sur P, alors il existe une unique décomposition de D en somme
disjointe d'ouverts $S_n (n \in \underline{N})$ telle que pour tout $n \in \underline{N}$ , $\underline{L}|P|S_n$ soit de
degré n. Cela ramène la détermination de l'ensemble des diviseurs rela\xout{xx}%
tifs $\geq 0$ au cas des diviseurs relatifs $\geq 0$ de degré n donné. (Ce
\margin{$/\!/$} \unsure{tèxte|texte} remplacera évidemment le ``abstraction faite'' plus haut). Ceci posé
les refexions précédentes établissent la

\uline{Proposition} 2.2. Sous les hypothèses précédentes, on a une correspondan%
ce biunivoque entre l'ensemble $\mathrm{Div}^{n}(\unsure{P}/S)$ des diviseurs relatifs $\geq 0$ de
degré n sur P, et l'ensemble des Modules quotients inversibles $\underline{M}$ de
$\underline{\mathrm{Sym}}^{n}(\underline{E})^{\hand{\vee}}$ , (ou ce qui revient au même, les sous-Modules inversibles loc
facteurs directs $\overset{\hand{\vee}}{\underline{M}}$ de $\underline{\mathrm{Sym}}^{n}(\underline{E})$ ) . Si D et $\underline{M}$ se correspondent, alors
$\underline{O}_P(D)$ est canoniquement isomorphe à $\underline{M}\otimes\underline{O}_P(n)$, et la section $s_D$ s'iden-%
tifie par cet isomorphisme à ce qu'on \unsure{devine}, compte tenu de $(*)$.

Notons que cette description est compatible avec tout changement de
base en S. Compte tenu de l'interprétation des Modules quotients inv de
$\underline{\mathrm{Sym}}^{n}(\underline{E})^{\hand{\vee}}$ comme les sections sur S de $P(\underline{\mathrm{Sym}}^{n}(\underline{E})^{\hand{\vee}})$, on trouve ceci:
\uline{Corollaire} 2.3. Le sous-foncteur \xout{xx} $\mathrm{Div}^{n}_{P/S}$ de $\mathrm{Div}^{+}_{\unsure{P}/S}$ est représentable
canoniquement par $P(\underline{\mathrm{Sym}}^{n}(\underline{E})^{\hand{\vee}})$.

Compte tenu des considérations de 2.1. , il s'ensuit:
\uline{Corollaire} 2.4. $\underline{\mathrm{Div}}^{+}_{\unsure{P}/S}$ est représentable canoniquement \xout{xxxxx} par
le S-préschéma somme des $P(\underline{\mathrm{Sym}}^{n}_{P/S}(\underline{E})^{\hand{\vee}})$ , $n \in \underline{N}$ .
2.5. Supposons maintenant donné une famille finie $(\underline{E}_i)_{i \in I}$ de Modules
localement libres sur S, d'où des $P_i = P(\underline{E}_i)$ et un P= produit fibré des
$P_i$ , appelé aussi \uline{fibré multiprojectif} sur S défini par $(\underline{E}_i)$ . Pour
simplifier les notations, désignons par $\underline{O}_i(n)$ l'image inverse sur P du

\mspage{107}
